\documentclass[11pt,a4paper,fontset=fandol]{ctexart}
\usepackage[margin=21mm]{geometry}
\usepackage{amsmath,amssymb,booktabs,array,tabularx,longtable}
\usepackage{enumitem,setspace,xcolor,hyperref,fancyhdr,needspace}
\hypersetup{colorlinks=true,linkcolor=blue!45!black,citecolor=blue!45!black,urlcolor=blue!45!black,pdftitle={Task II 共享计数方法与通用驻留端点}}
\setstretch{1.13}
\setlength{\emergencystretch}{3em}
\setlength{\headheight}{15pt}
\setlength{\tabcolsep}{5pt}
\renewcommand{\arraystretch}{1.18}
\setlist{nosep,leftmargin=2em}
\pagestyle{fancy}\fancyhf{}
\fancyhead[L]{\small Task II：共享计数方法}
\fancyhead[R]{\small Step 2，待审阅}
\fancyfoot[C]{\thepage}
\newcommand{\RI}{\mathrm{RI}}
\newcommand{\op}{\mathrm{op}}
\newcommand{\loc}{\mathrm{loc}}
\newcommand{\cn}[1]{c_{\mathrm N}\!\left(#1\right)}
\newcommand{\ck}[1]{c_{\mathrm K}\!\left(#1\right)}
\newcommand{\tN}{T_{\mathrm N}}
\newcommand{\tK}{T_{\mathrm K}}
\newcommand{\Byte}{\,\mathrm{Byte}}
\newcommand{\KiB}{\,\mathrm{KiB}}
\newcommand{\MiB}{\,\mathrm{MiB}}
\newcommand{\qs}{Q_{\mathrm S}}
\newcommand{\qr}{Q_{\mathrm R}}
\newcommand{\qso}{Q_{\mathrm S}^{\op}}
\newcommand{\qsl}{Q_{\mathrm S}^{\loc}}
\newcommand{\qro}{Q_{\mathrm R}^{\op}}
\newcommand{\qrl}{Q_{\mathrm R}^{\loc}}
\newcommand{\riop}{\RI^{\op}}
\newcommand{\riloc}{\RI^{\loc}}
\newcommand{\para}[1]{\par\smallskip\noindent\textbf{#1}\quad}
\title{\bfseries 共享计数方法与通用驻留端点\\[5pt]\large INT8 逻辑参考、128$\times$128 局部映射与四类矩阵模板}
\author{}\date{2026 年 9 月 24 日}
\begin{document}
\maketitle\vspace{-1.4em}
\noindent\textbf{状态：Step 1 已通过；Step 2 完成，待审阅。}
Step 1 审阅版本为 \texttt{e95de9e}，本轮从干净的 \texttt{d60bc32} 工作区继续。
本文固定共同计数规则，完成 Table II(a)；Table II(b) 仅给符号模板和合成检查，未代入六模型形成正式结果。

\section{分析前提、精度与算子交接}
\para{固定对象}
一次使用是一个输入向量求值。Table II(a) 的动态端点每次使用前装载完整矩阵；FFN 行的一次驻留则装载一次、服务 $B\in\{8,64,512\}$ 个向量。
Attention 的 $L\in\{1024,16384,131072\}$ 是本次追加后的可见长度。
六模型、revision 与选定的完整/全局 GQA 子层沿用 Step 1，均不乘主干层数。

\para{计量对象}
$\qs$ 是窗口内进入声明 streaming 输入边界并被服务的逻辑字节；$\qr$ 是窗口内建立、更新或重载可计算 resident 状态的逻辑写入字节\cite{conventions}。
当 $\qr>0$ 时 $\RI=\qs/\qr$；有求值且 $\qr=0$ 时 $\RI=\infty$。
空窗口的 $0/0$ 仅作递推检查，不给 RI 值。

\para{本轮新选的精度参考}
所有被计矩阵输入与 resident 数值均以每元素 1 Byte 的有符号 INT8 表示；Attention 系数使用其非负子域。符号推导仍分角色保留宽度：
\begin{center}\small
\begin{tabularx}{\textwidth}{@{}lXc@{}}\toprule
符号 & 数值角色 & 参考值（Byte/element）\\\midrule
$b_x,b_z$ & 投影/FFN 输入；门控乘积形成的 down 输入 & 1\\
$b_q,b_p$ & QK 求值的 query；softmax 后的 Attention 系数 & 1\\
$b_W$ & 通用矩阵权重 & 1\\
$b_g,b_u,b_d$ & FFN gate、up、down 三矩阵权重 & 1\\
$b_K,b_V$ & resident Key 与 Value & 1\\\bottomrule
\end{tabularx}\end{center}
这是与 Task I 的 INT8 逻辑合同对齐的研究参考，不改写已保存的原生 BF16/FP8 事实，也不构成六模型 INT8 精度验证。定标参数在数字接口预置，在所声明层/头/角色的窗口内保持固定；本参考计数不含定标元数据。

\para{阶段交接}
局部输出先在数字域做足够宽的部分和累加，再完成适用的归一化、RoPE、softmax、SiLU、逐元素乘法或 Value 缩放，最后按下一角色重定标/量化。
本次矩阵输出不直接加入 $\qs$；下次进入另一矩阵阶段时才计该阶段的输入。
位串行、ADC 转换、编码、verify、restore 等属于硬件服务，不重复增加逻辑 payload。
数字衔接与元数据虽不计入这里的矩阵数值流量，仍需真实实现的时间和资源。

\clearpage
\section{两种边界与共同驻留映射}
\para{主口径与对照口径}
上标 $\loc$ 表示所有 $128\times128$ tile 的\textbf{输入接收端累计}，用于与 Task I 的局部服务合同衔接；上标 $\op$ 表示\textbf{整矩阵阶段入口}，用于解释和对照。
同一阶段、同一源向量向并行分支分发时，$\op$ 只计一次上游输入；$\loc$ 对每个实际 tile 接收分别计数。
串行阶段的新输入再次计入，因此 FFN 的 $\op$ 仍含 down 输入，并非整个 FFN 外部 I/O 的唯一张量大小。

\para{分块与接收}
对 $W\in\mathbb Z^{N\times K}$，$N$ 是输出行、$K$ 是输入列。定义
\begin{equation}
\cn{x}=\left\lceil\frac{x}{\tN}\right\rceil,\qquad
\ck{x}=\left\lceil\frac{x}{\tK}\right\rceil,\qquad \tN=\tK=128.
\end{equation}
输出块 $a$、输入块 $b$ 的有效尺寸为 $n_a=\min(\tN,N-a\tN)$、$k_b=\min(\tK,K-b\tK)$。
每个 $(a,b)$ tile 常驻 $W_{a,b}$，一次向量求值接收 $k_b$ 个输入；同一输入分片可从上游缓冲广播，但不同 $a$ 的接收端分别计费。

\para{装载、求值与部分和}
动态驻留先把每个有效 tile 装载一次，随后按输出块、输入块的顺序枚举求值调用，并在数字累加器中归并同一输出块的部分和。
全部所需 tile 在声明驻留期内保留；静态权重装载在窗口之前。
该枚举是逻辑调用清单，不规定所有 tile 并行或某个最优时序。
参考有足够的逻辑服务单元容纳所声明状态，不引入容量不足的重载。
部分和不重新作为 CIM 输入，输入分块也不等于新增 bit-plane payload。

\para{可直接执行的通用式}
记通用输入与 resident 元素宽度为 $b_S,b_R$（Byte/element）。若窗口内有 $U$ 次向量使用、每个 resident tile 有 $A$ 次完整逻辑装载，则
\begin{align}
\qso&=UKb_S,&\qsl&=U\sum_a\sum_b k_b b_S=U\cn{N}Kb_S,\label{eq:generic-s}\\
\qro=\qrl&=A\sum_a\sum_b n_a k_b b_R=ANKb_R,&
C_{\mathrm{eval}}&=U\cn{N}\ck{K}.\label{eq:generic-r}
\end{align}
保留每种 $(n_a,k_b)$ 的有效形状和调用次数；逻辑 tile 装载项数是 $A\cn{N}\ck{K}$，不等于器件写事务数。
尾块只计有效元素，不据少输入字节推断完整服务时间同比缩短。

\para{合成例与结论}
仅为检查，令 $\tN=3,\tK=2,N=5,K=7,U=A=1$、每元素 1 Byte。
有效输出块为 $(3,2)$、输入块为 $(2,2,2,1)$，共 8 个 tile；得到
$(\qso,\qro)=(7,35)\Byte$、$(\qsl,\qrl)=(14,35)\Byte$。
改变为 $U=4,A=1$ 时，输入随使用次数变成 $28/56$ Byte，resident 写入仍为 35 Byte。

\para{与硬件能力的对应}
累计需求描述一个明确的 tile 集合，不能直接除以\emph{单 macro} 能力当成完成时间。
之后若进行硬件比较，必须对同一集合和精度声明资源归约、共享/互斥、尾块利用率及写入模式；本轮只给需求，不重估 Task I。

\clearpage
\section{Table II(a)：两种驻留端点}
\para{前提}
窗口只覆盖一个输入向量，故 $U=1$。Weight-static 取 $A=0$，权重在窗口外预驻留；Per-use reloaded 取 $A=1$，窗口内完整装载一次再求值。
后者不是完整训练 step，也不把多向量 GEMM 放进这次使用。
由式\eqref{eq:generic-s}--\eqref{eq:generic-r} 得到
\begin{center}\small
\begin{tabular}{@{}lll@{}}\toprule
驻留端点 & 整算子阶段入口 $(\qso,\qro,\riop)$ & 局部接收端 $(\qsl,\qrl,\riloc)$\\\midrule
Weight-static & $(Kb_S,0,\infty)$ & $(\cn{N}Kb_S,0,\infty)$\\[5pt]
Per-use reloaded & $\left(Kb_S,NKb_R,\dfrac{b_S}{Nb_R}\right)$ & $\left(\cn{N}Kb_S,NKb_R,\dfrac{\cn{N}b_S}{Nb_R}\right)$\\\bottomrule
\end{tabular}\end{center}
这两套三元组有不同输入边界；共享的 $\qr$ 只表示本参考没有复制权重或额外重载。

\para{精确数值}
下表取 $b_S=b_R=1$。写入列仅对应动态端点；静态端点的输入列相同、$\qr=0$、$\RI=\infty$。
\begin{center}\small
% Generated exact II(a) contrast values.
\begin{tabular}{@{}rrrrrr@{}}\toprule
$(N,K)$ & tile/调用数 & $\qso$ (Byte) & $\qsl$ (Byte) & 动态 $\qr$ (Byte) & 动态 $\riop$\\\midrule
$(128,128)$ & 1 & 128 & 128 & 16\,384 & $1/128$ \\
$(1024,1024)$ & 64 & 1\,024 & 8\,192 & 1\,048\,576 & $1/1024$ \\
$(4096,4096)$ & 1024 & 4\,096 & 131\,072 & 16\,777\,216 & $1/4096$ \\
$(1024,4096)$ & 256 & 4\,096 & 32\,768 & 4\,194\,304 & $1/1024$ \\
$(4096,1024)$ & 256 & 1\,024 & 32\,768 & 4\,194\,304 & $1/4096$ \\
\bottomrule\end{tabular}
\end{center}
\begin{equation}
\boxed{\quad \riloc_{\mathrm{reload}}=\frac{1}{128}=0.0078125
\quad\text{（五个规定形状均成立）}.\quad}\label{eq:IIa-ri}
\end{equation}

\para{为何相同}
五种形状的 $N$ 都是 128 的整数倍，因此 $\cn{N}=N/128$。
固定一次使用，局部输入为 $NK/128$ Byte，完整装载为 $NK$ Byte，比例即 $1/128$。
例如 $1024\times4096$ 和 $4096\times1024$ 的局部输入都为 32 KiB、写入都为 4 MiB；整算子输入却分别为 4 KiB 与 1 KiB。
规模不同不必导致比例不同；这里没有为制造差异更改边界。

\para{尾块说明}
对非整除 $N$，动态 $\riloc=\lceil N/128\rceil/N$（等字节宽度），不强行套用 $1/128$。
输入方向的尾块不增加有效逻辑总长 $K$，但调用数仍包含其独立服务。

\para{交付入口}
英文主表采用局部接收端口径，逐格列 $\qs,\qr,\RI$。
\path{table_IIa/data/results.json} 保存两套完整精确数据；
\path{table_IIa/tex/table_fragment.tex} 是生成片段，
\path{table_IIa/tex/table_IIa.tex} 是独立英文入口。
形状只定义逻辑任务，不重新定义五种物理 macro。

\clearpage
\section{QKV Projection：原生矩阵与统一端口规则}
\para{前提与布局选择}
输入为一个 $D$ 维 token 向量，投影权重预驻留。
本参考将语义矩阵 $Q$、可选的 Attention output gate $G$、$K$、$V$ 分别切成 tile bank，\textbf{不跨语义矩阵合并尾块}。
源码中的融合张量可按语义切片/重排；这一存储动作不复制权重，也不自动合并 CIM 接收端。
因此 fused QKV 与独立 Q/K/V 都遵循同一规则。

令 $H_q,H_{kv},d_{QK},d_V$ 为原生头数和维度，$D_g$ 为额外 output gate 宽度。无门控时 $D_g=0$；两款选定 Qwen 的 $D_g=H_qd_V$，其原生 q-projection 不能缩成纯 Q\cite{qwenimpl,models}。
\begin{center}\small
\begin{tabular}{@{}llll@{}}\toprule
矩阵 & 有效形状 $N\times K$ & 单独入口 $\qs^\op$ & 接收端 $\qs^\loc$\\\midrule
$W_Q$ & $(H_qd_{QK})\times D$ & $Db_x$ & $\cn{H_qd_{QK}}Db_x$\\
$W_G$（若有） & $D_g\times D$ & $Db_x$ & $\cn{D_g}Db_x$\\
$W_K$ & $(H_{kv}d_{QK})\times D$ & $Db_x$ & $\cn{H_{kv}d_{QK}}Db_x$\\
$W_V$ & $(H_{kv}d_V)\times D$ & $Db_x$ & $\cn{H_{kv}d_V}Db_x$\\\bottomrule
\end{tabular}\end{center}
每项 $\qr=0,\RI=\infty$，$W_G$ 不存在时整项省略。

\para{汇总与共享输入}
设实际非空矩阵集合为 $\mathcal J$，$m=|\mathcal J|$，各自输出宽度为 $n_j$。
同一个 token 输入在整算子阶段入口只进入一次，而每个 tile bank 接收端独立：
\begin{align}
Q_{\mathrm{S,QKV}}^{\op}
 &=\sum_{j\in\mathcal J}Db_x-(m-1)Db_x=Db_x,\\
Q_{\mathrm{S,QKV}}^{\loc}
 &=Db_x\sum_{j\in\mathcal J}\cn{n_j},&
Q_{\mathrm{R,QKV}}&=0,\\
C_{\mathrm{QKV}}&=\ck{D}\sum_{j\in\mathcal J}\cn{n_j},&
\RI_{\mathrm{QKV}}^{\op}=\RI_{\mathrm{QKV}}^{\loc}&=\infty.
\end{align}
上表的 $\op$ 分项是\emph{独立矩阵入口}，有同源重叠，汇总时已显式扣除；$\loc$ 分项的接收端不同，可以直接相加。
只按数值相等而去重不同 token/heads 的输入不属于本规则。

\para{合成检查}
设 $D=5,H_q=4,H_{kv}=2,d_{QK}=3,d_V=2,D_g=8$，测试 tile 为 $3\times2$。
四项输出宽度为 $12,8,6,4$，对应 $4,3,2,2$ 个输出块。
故 $\qso=5$ Byte、$\qsl=55$ Byte、调用数 $3(4+3+2+2)=33$，$\qr=0$。
没有 gate 时局部输入为 40 Byte；额外门控的投影需求已明确保留。

\para{算子交接与结论}
Q/K 的归一化和 RoPE、gate 的 sigmoid 与逐元素乘法在数字域处理；gate 本身不成为 KV 缓存，也不变成新的 query head。
本行仅含原生投影，不含 Attention 输出投影 $W_O$，不把 KV 写入混入预驻留投影权重的 $\qr$。
KV 建立/追加另由下文 Attention 窗口计量。各行是独立场景，不能不检查窗口就拼成端到端总量。

\clearpage
\section{FFN / MoE：一次装载后的实际向量复用}
\para{前提}
只分析一个 Dense FFN 或一个路由专家的完整 SwiGLU：
\begin{equation}
z=\operatorname{SiLU}(W_gx)\odot(W_ux),\qquad y=W_dz,
\end{equation}
其中 $W_g,W_u\in\mathbb Z^{F\times D}$、$W_d\in\mathbb Z^{D\times F}$\cite{glu,models}。
窗口内三矩阵各装载一次，服务 $B$ 个实际输入向量。三个 tile bank 同期保留，无专家替换的额外重载；窗口结束后替换不在本窗口重复收费。

\para{分项需求}
$z$ 是数字域 SiLU 和逐元素乘积完成后，重定标成宽度 $b_z$ 的新矩阵输入。
\begin{center}\small
\begin{tabular}{@{}llll@{}}\toprule
矩阵 & 单独入口 $\qs^\op$ & 接收端 $\qs^\loc$ & $\qr$（两边界一致）\\\midrule
$W_g$ & $BDb_x$ & $B\cn{F}Db_x$ & $FDb_g$\\
$W_u$ & $BDb_x$ & $B\cn{F}Db_x$ & $FDb_u$\\
$W_d$ & $BFb_z$ & $B\cn{D}Fb_z$ & $DFb_d$\\\bottomrule
\end{tabular}\end{center}
每项 RI 按本行对应两列求商，例如
\begin{equation}
\RI_g^\op=\frac{Bb_x}{Fb_g},\quad
\RI_g^\loc=\frac{B\cn{F}b_x}{Fb_g},\qquad
\RI_d^\loc=\frac{B\cn{D}b_z}{Db_d}.
\end{equation}
up 将 $b_g$ 换为 $b_u$，down 的整算子式为 $Bb_z/(Db_d)$。

\para{汇总公式}
gate/up 使用同一份 $x$，其上游输入可共享；down 使用另一个 $z$。
两组阶段入口之间不把 $W_gx,W_ux$ 输出再各加一份，只在 down 真正接收 $z$ 时计入：
\begin{align}
Q_{\mathrm{S,FFN}}^{\op}&=B(Db_x+Fb_z),\label{eq:ffn-o}\\
Q_{\mathrm{S,FFN}}^{\loc}&=B\left(2\cn{F}Db_x+\cn{D}Fb_z\right),\label{eq:ffn-l}\\
Q_{\mathrm{R,FFN}}&=DF(b_g+b_u+b_d),\label{eq:ffn-r}\\
C_{\mathrm{FFN}}&=B\left(2\cn{F}\ck{D}+\cn{D}\ck{F}\right).
\end{align}
式\eqref{eq:ffn-o} 从独立分项之和扣除一次共享的 $BDb_x$；式\eqref{eq:ffn-l} 保留两个 bank 的接收。
两种汇总 RI 都由对应的累计 $\qs/\qr$ 得到，不能平均分项 RI。

\para{合成例}
令 $D=5,F=7,B=4$，测试 tile 为 $3\times2$，所有字节宽度为 1。
局部分项输入为 gate/up/down $=(60,60,56)$ Byte，三项写入各 35 Byte。
于是局部汇总为 $(176,105,176/105)$，整算子汇总为 $(48,105,16/35)$。
增加 $B$ 只增加使用次数和输入接收；本驻留期的完整装载仍只有一次。

\para{对象范围}
MoE 的 top-k、全部路由专家、共享专家和其它层均不作本行隐式倍数。
路由矩阵不属于所选专家的 gate/up/down；也不设路由均匀、命中率或额外切换假设。
$B$ 是同一份专家权重真正服务的向量总数，不是批次数或全局 batch size。

\clearpage
\section{Attention：两种 resident 方向与单前缀规则}
\para{头参数与 GQA}
保留独立的 $H_q,H_{kv},d_{QK},d_V$，令 $g=H_q/H_{kv}$ 为整数。
每个 KV head 只持有一份 K/V，组内 $g$ 个 query heads 依次对同一 resident 状态求值\cite{gqa}。
每个 query head 有自己的 query 和 Attention 系数，streaming 按各次求值计数；源码中的 \texttt{repeat\_kv} 不形成额外持久副本。

\para{固定状态坐标与方向}
对 KV head $h$、可见前缀 $i$，把每个 token 的 Key 和 Value 分别排列为
\begin{equation}
\mathcal K_i^{(h)}=
\begin{bmatrix}k_1^{(h)\mathsf T}\\\vdots\\k_i^{(h)\mathsf T}\end{bmatrix}
\in\mathbb R^{i\times d_{QK}},\qquad
\mathcal V_i^{(h)}=
\begin{bmatrix}v_1^{(h)}&\cdots&v_i^{(h)}\end{bmatrix}
\in\mathbb R^{d_V\times i}.
\end{equation}
列向量写法下，$s_i^{(q)}=\mathcal K_i^{(h)}q_i^{(q)}$，
$o_i^{(q)}=\mathcal V_i^{(h)}p_i^{(q)}$，与常见行向量 $QK^{\mathsf T}$ 和 $PV$ 等价\cite{attention}。
\textbf{K 的序列方向是输出行，V 的序列方向是输入列。}
因此序列增大导致 QK 的 query 向更多输出块重放，而 AV 的系数向量本身变长。

\para{因果前缀的求值组织}
Prefill 从空本段状态开始：对 $i=1,\ldots,L$，先将 token $i$ 的 K/V 各写入一次，再对全部 $H_q$ 个 query heads 求值可见前缀 $1{:}i$；计算所有有效 scores 后做因果 softmax，再求值 AV。
这是完整的逐前缀逻辑参考，不声称是推理引擎的最快 Prefill 调度。
旧 token 不移动、不压缩、不重新装载；未来未写入位置不参加本次求值。

\para{一个前缀的输入与调用}
同一 token 下所有 query heads 合计：
\begin{align}
Q_{\mathrm S,QK}^{\op}(i)&=H_qd_{QK}b_q,&
Q_{\mathrm S,QK}^{\loc}(i)&=H_qd_{QK}b_q\cn{i},\label{eq:prefix-qk}\\
Q_{\mathrm S,AV}^{\op}(i)&=H_q i b_p,&
Q_{\mathrm S,AV}^{\loc}(i)&=H_q i b_p\cn{d_V},\label{eq:one-prefix}\\
C_{QK}(i)&=H_q\cn{i}\ck{d_{QK}},&
C_{AV}(i)&=H_q\cn{d_V}\ck{i}.\label{eq:prefix-calls}
\end{align}
$\qs$ 只包含 query 和下一次求值入口的系数 $p$。scores 是第一次输出，不另加到 $\qs$。
softmax 需要的数字保存与处理、RoPE/归一化及 Value 缩放均在矩阵接口交接。

\para{追加的形状与写入}
每个 KV head 追加 K 的 $1\times d_{QK}$ 行、V 的 $d_V\times1$ 列：
\begin{equation}
\Delta Q_{\mathrm R,K}=H_{kv}d_{QK}b_K,\qquad
\Delta Q_{\mathrm R,V}=H_{kv}d_Vb_V.\label{eq:append}
\end{equation}
K 行分散到 $\ck{d_{QK}}$ 个 $1\times k_b$ 写入片；V 列分散到 $\cn{d_V}$ 个 $n_a\times1$ 写入片。
在预分配的绝对 token 坐标上增长有效形状，跨 tile 边界时只写新元素，不把所有小更新补成整 tile 重写。
这些逻辑形状尚不保证任何介质的物理 append 效率。

\clearpage
\section{Prefill 与 Decode 的精确闭式模板}
\para{因果求和}
从相同前缀规则精确累加，定义
\begin{align}
P(L)&=\sum_{i=1}^{L}i=\frac{L(L+1)}2,\\
S_T(L)&=\sum_{i=1}^{L}\left\lceil\frac iT\right\rceil
       =\frac{Tm(m+1)}2+(m+1)r,\quad L=mT+r,\quad 0\le r<T.\label{eq:st}
\end{align}
每个完整的 $T$ 项区间贡献 $T,2T,\ldots,mT$，余下 $r$ 项各贡献 $m+1$。
这里不以 $L^2/2$ 替代 $P(L)$，也不以连续近似代替 ceiling 求和。

\para{Prefill：从空 KV 到长度 $L$}
两次矩阵任务的分项与汇总为
\begin{align}
Q_{\mathrm S,QK}^{\op}&=H_qLd_{QK}b_q,&
Q_{\mathrm S,QK}^{\loc}&=H_qd_{QK}b_q S_{\tN}(L),\\
Q_{\mathrm S,AV}^{\op}&=H_qb_p P(L),&
Q_{\mathrm S,AV}^{\loc}&=H_q\cn{d_V}b_p P(L),\\
Q_{\mathrm R,K}&=LH_{kv}d_{QK}b_K,&
Q_{\mathrm R,V}&=LH_{kv}d_Vb_V,\\
Q_{\mathrm{S,pre}}^{\op}&=H_q\left[Ld_{QK}b_q+b_pP(L)\right],\\
Q_{\mathrm{S,pre}}^{\loc}&=H_q\left[d_{QK}b_q S_{\tN}(L)+\cn{d_V}b_pP(L)\right],\\
Q_{\mathrm{R,pre}}&=LH_{kv}(d_{QK}b_K+d_Vb_V),\label{eq:pre-r}\\
C_{\mathrm{pre}}&=H_q\left[\ck{d_{QK}}S_{\tN}(L)+\cn{d_V}S_{\tK}(L)\right].
\end{align}
每份 KV 会被多个后续 query 向量读取；“一次 Attention 任务”并不等于每份 KV 只服务一次。

\para{Decode：从 $L-1$ 追加到可见 $L$}
$\qs$ 与调用数分别取式\eqref{eq:prefix-qk}--\eqref{eq:prefix-calls} 中的 $i=L$，写入取式\eqref{eq:append}：
\begin{align}
Q_{\mathrm{S,dec}}^{\op}&=H_q(d_{QK}b_q+Lb_p),\\
Q_{\mathrm{S,dec}}^{\loc}&=H_q\left[d_{QK}b_q\cn{L}+Lb_p\cn{d_V}\right],\\
Q_{\mathrm{R,dec}}&=H_{kv}(d_{QK}b_K+d_Vb_V),\label{eq:dec-r}\\
C_{\mathrm{dec}}&=H_q\left[\cn{L}\ck{d_{QK}}+\cn{d_V}\ck{L}\right].
\end{align}
写入只计新 token，已有 $L-1$ 个 token 是窗口初态。对每个分项和汇总，$\RI$ 都是其累计 $\qs/\qr$；$\op$、$\loc$ 不能混用，分项 RI 不做算术平均。

\clearpage
\section{合成例、独立核验与适用条件}
\para{不等宽 GQA 例}
取 $H_q=4,H_{kv}=2,d_{QK}=3,d_V=5,L=7$，测试 tile 为 $3\times2$，所有数值每元素 1 Byte。
$P(7)=28$、$S_3(7)=12$、$S_2(7)=16$。
\begin{center}\small
\begin{tabular}{@{}llrrrr@{}}\toprule
窗口 & 分项 & $\qs^\op$ (Byte) & $\qs^\loc$ (Byte) & $\qr$ (Byte) & tile 调用\\\midrule
Prefill & QK & 84 & 144 & 42 & 96\\
 & AV & 112 & 224 & 70 & 128\\
 & 合计 & 196 & 368 & 112 & 224\\\midrule
Decode & QK & 12 & 36 & 6 & 24\\
 & AV & 28 & 56 & 10 & 32\\
 & 合计 & 40 & 92 & 16 & 56\\\bottomrule
\end{tabular}\end{center}
Prefill 的局部 RI 为 $368/112=23/7$，Decode 为 $92/16=23/4$。
Prefill 分项 RI 的算术平均为 $\tfrac12(144/42+224/70)=116/35$，并不等于 $23/7$；应先相加字节再求商。
最终只有 $2(7\times3+5\times7)=112$ 个 resident 元素，不随 $g=2$ 复制成 224 个。

\para{窗口衔接}
相同组织下，Prefill 到 6 的局部输入/写入/调用为 $(276,96,168)$。
追加第 7 个 token 的差量为 $(92,16,56)$，与 Decode 可见 7 的结果逐项一致。
一般有
\begin{equation}
\mathcal D_{\mathrm{pre}}(L)-\mathcal D_{\mathrm{pre}}(L-1)
=\mathcal D_{\mathrm{dec}}(L),\quad
\mathcal D\in\{\qs^\op,\qs^\loc,\qr,C_{\mathrm{eval}}\}.
\end{equation}
该差分适用于需求与调用数，\textbf{不适用于 RI 相减}。
从空 KV 逐 token 执行并累计 Decode 事件，即得到这里的 Prefill。

\para{独立检查方式}
\path{shared/scripts/counting.py} 实现闭式公式；
\path{check_counting.py} 另写分块求和及显式事件枚举，不复用闭式求值函数作为预期值。
显式枚举按输入身份、tile 接收端、求值编号和 resident 元素写入建立事件集合，逐项验证是否引用已建立的 K/V。
检查包含矩形、两方向非整除尾块、静态、一次装载多次使用、每次重载、额外 Q gate、分支共享、GQA、不等宽 QK/V、$L=1$ 与跨 128 边界，以及精确 Prefill/Decode 衔接。
Table II(a) 全部 10 格另由独立 tile 求和复算。

\para{集中的适用条件}
这是一套单一、可复算的逻辑映射：足够驻留空间、固定 tile 坐标、独立语义 bank、数字域累加与非矩阵交接。
角色字节宽度可以变化，但必须保持相同计量与对应能力配置；若需要多个 Task I 服务组合才能实现其它精度，应重新声明映射。
K 行追加与 V 列追加的写入形状已经保留；页/块擦写、内部读改写及最小事务约束不在这里凭空扩大逻辑 $\qr$，未来介质映射须据真实写服务另做代价/可行性判断。
本轮不从总需求推算延迟、token/s、端到端性能或硬件适合度。

\clearpage
\section{复现入口、证据与停止点}
\para{交付与生成关系}
\begin{itemize}
\item 本文：\path{shared/tex/counting_method.zh.tex}；独立 PDF 位于 \path{shared/output/counting_method.zh.pdf}。
\item 机器约定：\path{shared/data/conventions.json}，记录角色精度、双边界、窗口、语义 bank、状态方向及公式输入。
\item 合成检查：\path{shared/data/synthetic_checks.json}，原始定义与结果可逐项对应本稿例子。
\item Table II(a)：\path{table_IIa/data/results.json}、独立表格片段及英文 PDF；所有数值由 \path{generate_IIa.py} 生成。
\item 原始模型数据与资料仍由 \path{data/models.json}、\path{data/sources.json} 和 \path{literature/} 提供，不改写 Step 1 原生参数。
\end{itemize}
所有路径在未特别说明时相对本任务目录。执行
\begin{quote}\small
\texttt{TASK2\_PYTHON=/opt/anaconda3/bin/python \textbackslash}\\
\texttt{sh }\path{tasks/task2_table_II_workloads/shared/scripts/build.sh}
\end{quote}
先复算生成数据的一致性，再分别编译中文方法与英文表格；检查默认不覆盖数据，刷新需显式使用 \texttt{--emit}。
页面渲染与检查记录保存在各自 \path{tmp/pdfs/} 和 \path{shared/data/pdf_qa.json}。

\para{审阅重点}
本轮建议确认三项共同选择：INT8 的分角色接口参考；独立语义 bank（不跨分支合并尾块）；从空 KV 逐前缀建立并求值的因果组织。
这些是已落实的主参考，不要求先枚举多种优化调度。
\textbf{Step 3 与六模型 Table II(b) 正式数值仍未启动。}

\begin{thebibliography}{9}\small
\bibitem{conventions} 本仓库，\emph{CIM Roofline 计量约定}，\path{docs/MODEL_CONVENTIONS.md}；Task I 共享外围参考与十例审阅，\path{tasks/task1_table_I_NVM/analysis/}。本轮保持不变。
\bibitem{gqa} J. Ainslie et al., \emph{GQA: Training Generalized Multi-Query Transformer Models from Multi-Head Checkpoints}, arXiv:2305.13245v3, 2023，p.2，§2.2/Fig.2。本地：\path{literature/00_shared/raw/2305.13245v3.pdf}。
\bibitem{attention} A. Vaswani et al., \emph{Attention Is All You Need}, arXiv:1706.03762v7, pp.4--5，§3.2。本地：\path{literature/00_shared/raw/1706.03762v7.pdf}。
\bibitem{glu} N. Shazeer, \emph{GLU Variants Improve Transformer}, arXiv:2002.05202v1, pp.1--2，§2。本地：\path{literature/00_shared/raw/2002.05202v1.pdf}。
\bibitem{qwenimpl} Transformers 固定 commit \texttt{2c4914f}，\path{modeling_qwen3_5.py} 第 749--821 行，\path{modeling_qwen3_5_moe.py} 第 751 行起。原件位于 \path{literature/00_shared/raw/transformers/}；Q 与输出 gate 的投影/切分见 Attention 构造与 forward。
\bibitem{models} Step 1 固定六模型结构卡与 \path{data/models.json}。Ling/MiMo 原生实现位于各自资料包，MiMo 的 $d_{QK}\ne d_V$、Value 缩放及 K/V cache 更新均已定位。上游格式用于结构依据，本稿的 INT8 与端口映射是本文另选参考。
\end{thebibliography}
\end{document}
